\BourbakiExerciseGroup

\begin{enumerate}
\def\labelenumi{\arabic{enumi})}
\tightlist
\item
  设 \(f\) 是拓扑空间 \(\mathbf{X}\) 到拓扑空间 \(\mathbf{X}'\) 中的一个映射。证明以下命题等价：
\end{enumerate}

\begin{enumerate}
\def\labelenumi{\alph{enumi})}
\item
  \(f\) 在 \(\mathbf{X}\) 上是连续的。
\item
  对每个子集 \(\mathbf{A}'\)，有 \(\mathbf{X}', f^1 (\mathring{\mathbf{A}}') \subset (f^1 (\mathbf{A}'))^0\)。
\item
  对每个子集 \(\mathbf{A}'\)，有 \(\mathbf{X}', \overline{f}^1 (\mathbf{A}') \subset \overline{f}^1 (\overline{\mathbf{A}}')\)。
\end{enumerate}

用一个例子证明：若 \(f\) 连续，集合 \(\overline{f}^1 (\mathbf{A}')\) 与 \(\overline{f}^1 (\overline{\mathbf{A}}')\) 仍可以不同。

\begin{enumerate}
\def\labelenumi{\arabic{enumi})}
\setcounter{enumi}{1}
\item
  设 X, X' 是两个有序集，各赋予右拓扑（§ 1，习题 2）。证明映射 \(f: \mathbf{X} \to \mathbf{X}'\) 连续当且仅当它是保序的。
\item
  设 X, X' 是两个拓扑空间，f 是 X 到 X' 上的一个双射；则 f 是 X 到 \(X'\) 上的同胚的必要充分条件是：\(X'\) 的拓扑是使 f 连续的最细拓扑。
\item
  在有序集 X 上，左拓扑与右拓扑（§ 1，习题 2）的上确界是离散拓扑；若 X 是有向集（向右或向左有向），则这两个拓扑的下确界是 X 上最粗的拓扑。
\item
  在有序集 X 上，用 \(\mathcal{C}_{0}(\mathbf{X})\)（相应地，\(\mathcal{C}_{+}(\mathbf{X})\)，\(\mathcal{C}_{-}(\mathbf{X})\)）表示由所有开区间（相应地，右半开区间所成的集、左半开区间所成的集）生成的拓扑，这些区间可以有界也可以无界。
\end{enumerate}

\begin{enumerate}
\def\labelenumi{\alph{enumi})}
\item
  证明：若 X 是线性序的，则开（相应地，右半开、左半开）区间构成 \(\mathcal{C}_{0}(\mathbf{X})\)（相应地，\(\mathcal{C}_{+}(\mathbf{X})\)，\(\mathcal{C}_{-}(\mathbf{X})\)）的一个基；且闭区间在这三个拓扑的每一个中都是闭集。
\item
  在线性序集 X 上，证明拓扑 \(\mathcal{C}_{+}(\mathbf{X})\) 比右拓扑（§ 1，习题 2）细，拓扑 \(\mathcal{C}_{-}(\mathbf{X})\) 比左拓扑细，且拓扑 \(\mathcal{C}_{0}(\mathbf{X})\) 是 \(\mathcal{C}_{+}(\mathbf{X})\) 与 \(\mathcal{C}_{-}(\mathbf{X})\) 的交。
\item
  设 X 是一个线性序集。证明 \(\mathfrak{C}_{0}(\mathbf{X})\) 具有可数基的必要充分条件是：存在 X 的一个可数子集 D，使得对每个 \(x \in X\) 以及每个含有 x 的区间 {]}a, b{[}，存在 \(\alpha \in D\) 与 \(\beta \in D\) 使得
\end{enumerate}

\[
a \leqslant \alpha <   x <   \beta \leqslant b
\]

{[}参见第四章，§ 2，习题 11 c){]}。

\begin{enumerate}
\def\labelenumi{\alph{enumi})}
\setcounter{enumi}{3}
\tightlist
\item
  证明：若 X 是良序的（《集合论》，第三章，§ 2），则拓扑 \(\mathcal{G}_{+}(\mathbf{X})\) 是离散拓扑，且拓扑 \(\mathcal{G}_{0}(\mathbf{X})\) 与 \(\mathcal{G}_{-}(\mathbf{X})\) 相同。
\end{enumerate}

¶ 6) 集合 X 上的一个拓扑称为拟极大的，如果它在使 X 没有孤立点的诸拓扑所成的集中是极大的。

\begin{enumerate}
\def\labelenumi{\alph{enumi})}
\item
  证明 X 上的拓扑 \(\mathcal{G}\) 的以下性质等价：\(\alpha\) \(\mathcal{G}\) 是拟极大的；\(\beta\) 若 X 带有拓扑 \(\mathcal{G}\)，则 X 没有孤立点，且 X 的每个没有孤立点的子集都是开的。
\item
  设 X 是一个 Kolmogoroff 空间（§ 1，习题 2），其中每个非空开集都是无限的。证明 X 上存在一个比给定拓扑更细的拟极大拓扑（证明 X 上使所有非空开集都是无限的诸拓扑所成的集是归纳的，并用 § 1 的习题 2 d)）。
\end{enumerate}

\begin{enumerate}
\def\labelenumi{\arabic{enumi})}
\setcounter{enumi}{6}
\tightlist
\item
  若 X 是任一集合，用 \(\mathfrak{P}_{0}(\mathbf{X})\) 表示 X 的非空子集所成的集。
\end{enumerate}

\begin{enumerate}
\def\labelenumi{\alph{enumi})}
\item
  设 X 是一个非空拓扑空间。用 \(\mathcal{G}_{\Omega}\)（相应地，\(\mathcal{G}_{\Phi}\)）表示 \(\mathfrak{P}_{0}(\mathbf{X})\) 上使下列条件成立的最粗拓扑：对 X 的每个非空开（相应地，闭）子集 A，\(\mathfrak{P}_{0}(\mathbf{A})\) 在 \(\mathfrak{P}_{0}(\mathbf{X})\) 中是开的（相应地，闭的）。证明这两个拓扑一般不可比较，且映射 \(x \to \{x\}\) 是 X 到 \(\mathfrak{P}_{0}(\mathbf{X})\) 的一个子空间上的同胚，只要后一集合带有拓扑 \(\mathcal{G}_{\Omega}, \mathcal{G}_{\Phi}\) 之一。
\item
  设 D 是 X 的一个稠密子集。证明 D 的有限非空子集所成的集 \(\mathfrak{E}(D)\) 在 \(\mathfrak{P}_{0}(X)\) 中对于拓扑 \(C_{\Omega}, C_{\Phi}\) 中的每一个都是稠密的。
\item
  若 \(\mathfrak{P}_{0}(\mathbf{X})\) 按包含关系排序，证明拓扑 \(\mathcal{C}_{\Omega}\) 比 \(\mathfrak{P}_{0}(\mathbf{X})\) 上的左拓扑（\(\S\) 1，习题 2）粗；\(\mathbf{A} \in \mathfrak{P}_{0}(\mathbf{X})\) 在拓扑 \(\mathcal{C}_{\Omega}\) 中是孤立点，当且仅当 \(\mathbf{A} = \{x\}\)（其中 \(x\) 是 \(\mathbf{X}\) 的一个孤立点）；且 \(\mathbf{A} \in \mathfrak{P}_{0}(\mathbf{X})\) 在拓扑 \(\mathcal{C}_{\Phi}\) 中是孤立点，当且仅当 \(\mathbf{X}\) 是有限离散空间且 \(\mathbf{A} = \mathbf{X}\)。
\item
  证明：若 \(\mathfrak{P}_0(\mathbf{X})\) 带有拓扑 \(\mathcal{C}_{\Omega}\)（相应地，\(\mathcal{C}_{\Phi}\)）且
\end{enumerate}

\[
\mathfrak {P} _ {\mathbf {0}} (\mathbf {X}) \times \mathfrak {P} _ {\mathbf {0}} (\mathbf {X})
\]

带有 \(\mathcal{C}_{\Omega}\)（相应地，\(\mathcal{C}_{\Phi}\)）与其自身的积拓扑，则映射 \((\mathbf{M}, \mathbf{N}) \to \mathbf{M} \cup \mathbf{N}\)（\(\mathfrak{P}_0(\mathbf{X}) \times \mathfrak{P}_0(\mathbf{X})\) 到 \(\mathfrak{P}_0(\mathbf{X})\) 中的映射）是连续的。

\begin{enumerate}
\def\labelenumi{\alph{enumi})}
\setcounter{enumi}{4}
\tightlist
\item
  证明：若 \(\mathfrak{P}_{0}(\mathbf{X})\) 带有拓扑 \(\mathfrak{T}_{\Phi}\)，则映射 \(\mathbf{M} \to \overline{\mathbf{M}}\) 是连续的。
\end{enumerate}

\(f)\) 设 \(\mathbf{X}\), \(\mathbf{Y}\) 是两个拓扑空间，\(f: \mathbf{X} \to \mathbf{Y}\) 是连续映射。证明：若 \(\mathfrak{P}_{0}(\mathbf{X})\) 与 \(\mathfrak{P}_{0}(\mathbf{Y})\) 都带有拓扑 \(\mathcal{G}_{\Omega}\)，且都带有拓扑 \(\mathcal{G}_{\Phi}\)，则映射 \(\mathbf{M} \to f(\mathbf{M})\)（\(\mathfrak{P}_{0}(\mathbf{X})\) 到 \(\mathfrak{P}_{0}(\mathbf{Y})\) 中的映射）是连续的。

\begin{enumerate}
\def\labelenumi{\alph{enumi})}
\setcounter{enumi}{6}
\tightlist
\item
  设 X, Z 是两个拓扑空间。证明：映射 \(f: Z \to \mathfrak{P}_{0}(X)\) 当 \(\mathfrak{P}_{0}(X)\) 带有拓扑 \(C_{\Omega}\)（相应地，\(C_{\Phi}\)）时连续的必要充分条件是：对 X 的每个闭（相应地，开）子集 A，由 \(z \in Z\)（满足 \(f(z) \cap A \neq \emptyset\) 者）全体所成的集在 Z 中是闭的（相应地，开的）。
\end{enumerate}

\begin{enumerate}
\def\labelenumi{\arabic{enumi})}
\setcounter{enumi}{7}
\tightlist
\item
  设 X 是一个集合，c 是一个基数，它或是无限或等于 1。由 X 与那些满足 card(M) \textless{} c 的 M ⊆ X 组成的 X 的子集所成的集，是 X 关于某个拓扑 \(\mathfrak{G}_{\mathfrak{c}}\) 的闭子集所成的集；\(\mathfrak{G}_{\mathfrak{1}}\) 是 X 上最粗的拓扑，而 \(\mathfrak{G}_{\mathfrak{c}}\) 在 c \textgreater{} card (X) 时是离散拓扑。X 的每个置换都是把 X 赋予拓扑 \(\mathfrak{G}_{\mathfrak{c}}\) 所得拓扑空间的一个自同构。
\end{enumerate}

证明反之，任何拓扑 \(\mathcal{C}\)（在 \(\mathbf{X}\) 上）若具有性质：\(\mathbf{X}\) 的每个置换关于 \(\mathcal{C}\) 都连续，则它必然是诸拓扑 \(\mathcal{C}_{\mathfrak{c}}\) 之一。{[}设 \(\mathfrak{c}\) 是使 card \((\mathbf{F}) < \mathfrak{c}\)（对每个闭子集 \(\mathbf{F} \neq \mathbf{X}\)，在拓扑 \(\mathcal{C}\) 中）成立的最小基数；注意当 \(\mathbf{X}\) 无限且 \(\mathfrak{c} >\mathrm{card}(\mathbf{X})\) 时，\(\mathcal{C}\) 必然是离散拓扑。{]}

\begin{enumerate}
\def\labelenumi{\arabic{enumi})}
\setcounter{enumi}{8}
\tightlist
\item
  \begin{enumerate}
  \def\labelenumii{\alph{enumii})}
  \tightlist
  \item
    在第 3 目命题 4 的假设下，证明 X 上使诸 \(f_{t}\) 连续的最粗拓扑也是 X 上具有以下性质的最细拓扑 \(\mathfrak{G}\)：从拓扑空间 Z 到 X 中的每个映射 \(g\)（使诸 \(f_{t} \circ g\) 连续者）都（关于 \(\mathfrak{G}\)）连续。
  \end{enumerate}
\end{enumerate}

\begin{enumerate}
\def\labelenumi{\alph{enumi})}
\setcounter{enumi}{1}
\tightlist
\item
  在第 4 目命题 6 的假设下，证明 X 上使诸 \(f_{i}\) 连续的最细拓扑也是 X 上具有以下性质的最粗拓扑 \(\mathcal{G}\)：从 X 到拓扑空间 Z 中的每个映射 \(g\)（使诸 \(g \circ f_{i}\) 连续者）都（关于 \(\mathcal{G}\)）连续。
\end{enumerate}

\begin{enumerate}
\def\labelenumi{\arabic{enumi}.}
\setcounter{enumi}{9}
\tightlist
\item
  在集合 X 上，设 M 是由没有孤立点的拓扑组成的一个有向集。证明 M 在 X 上所有拓扑所成的集中的上确界是一个没有孤立点的拓扑。
\end{enumerate}

¶ 11. 拓扑空间称为可解的，如果它是两个稠密子集的不相交并。

\begin{enumerate}
\def\labelenumi{\alph{enumi})}
\item
  若拓扑空间 X 可解，证明 X 没有孤立点，且 X 的每个开子空间都可解。再证明反之，若 X 是拓扑空间，B 是 X 的非空可解开子空间所成的一个集，使得 X 的每个非空开子集都含有属于 B 的一个集合，则 X 可解（考虑属于 B 的两两不相交开集的一个极大集）。
\item
  若 Kolmogoroff 空间 X 可解，则 X 的每个非空开子集都是无限的{[}用 §1 的习题 2d){]}。
\item
  设 X 是这样的拓扑空间：1) X 的非空开子集的最小基数 b 是无限的；2) X 的拓扑有一个基 B 使得 card (B) ≤ b。证明 X 可解。（用超限归纳法构造 X 的两个不相交稠密子集。）用《集合论》，第三章，§ 6，习题 24 a) 的方法。特别地，有理直线是可解的。
\item
  没有孤立点的拓扑空间 X 称为等势的，如果对 X 的每个非空开子集有 card (U) = card (X)。证明在一个没有孤立点的空间 X 中，X 的每个非空开集都含有一个非空的等势开子空间。由此推出，若 X 的每个等势开子空间都可解，则 X 可解。
\end{enumerate}

